\begin{table}[h]
\footnotesize
\centering
\begin{tabular}{|p{4cm}|p{12.5cm}|}
\hline
\textbf{Título:} & SISTEMA DE VISIÓN ARTIFICIAL PARA LA DETECCIÓN Y CORRECCIÓN DE POSTURAS EN EJERCICIOS FÍSICOS \\ \hline

\textbf{Palabras clave:} & Visión artificial, estimación de pose humana, corrección postural, ejercicio de fuerza, análisis biomecánico. \\ \hline

\textbf{Resumen:} &
\begin{minipage}[t]{\linewidth}
\justifying
Este proyecto presenta el desarrollo de un sistema de visión artificial orientado a la detección y corrección de posturas durante la ejecución de ejercicios físicos de fuerza. Mediante el uso de modelos avanzados de estimación de pose humana, el sistema analiza en tiempo real la posición de las articulaciones del usuario para identificar desviaciones respecto a patrones biomecánicos óptimos. Se enfoca en cuatro ejercicios fundamentales del fisicoculturismo: sentadilla, peso muerto, curl de bíceps y press de banca, considerando criterios técnicos como ángulos articulares, simetría corporal y alineación postural. El sistema proporciona retroalimentación visual y textual personalizada para mejorar la técnica del usuario, prevenir lesiones y optimizar el rendimiento muscular. Esta solución busca ofrecer una alternativa accesible y automatizada al acompañamiento profesional en contextos de entrenamiento físico.
\end{minipage} \\ \hline

\textbf{Área temática:} & Procesamiento de Imágenes \\ \hline

\textbf{Modalidad de trabajo de grado:} &
$\boxtimes$ Profesional \quad $\Box$ Revisión crítica \quad $\Box$ Investigación e innovación \\
& $\Box$ En la industria/empresa \quad $\Box$ Creación de empresa \quad $\Box$ Práctico social \\ \hline

\textbf{Duración del trabajo de grado:} & 10 Meses \\ \hline
\textbf{Entidades participantes:} & Universidad del Valle \\ \hline

\textbf{Costo presupuestado del trabajo:} &
\makecell[tl]{%
Aportes en especie: COP\$ 36'160.000\\
Aportes en efectivo: COP\$ 0\\
Costo global: COP\$ 36'160.000
} \\ \hline

\begin{comment}
 \textbf{Áreas académicas EIEE:} &
\makecell[tl]{%
$\Box$ Arquitecturas Digitales\\
$\Box$ Bionanoelectrónica\\
$\Box$ Conversión de energía\\
$\Box$ Informática industrial\\
$\Box$ Sistema de control y accionamientos\\
$\Box$ Sistemas de potencia\\
$\Box$ Telecomunicaciones
} \\ \hline   
\end{comment}



\textbf{Grupos de investigación:} &
\makecell[tl]{%
$\Box$ Alta tensión (GRALTA)\\
$\Box$ Arquitecturas digitales y microelectrónica (GADYM)\\
$\Box$ Bionanoelectrónica\\
$\Box$ Control industrial (GICI)\\
$\Box$ Conversión de energía (CONVERGIA)\\
$\boxtimes$ Percepción y sistemas inteligentes (PSI)\\
$\Box$ Sistemas de telecomunicaciones (SISTEL-UV)\\
$\Box$ Sistema de hidroeléctricos de generación (SIIG)\\
$\Box$ Otro:
} \\ \hline



\multicolumn{2}{|c|}{\textbf{Estudiante}} \\ \hline
\textbf{Nombre:} & Susana Valencia Bravo \hfill \textbf{Código:} 202179092 \\ \hline
\textbf{Correo electrónico:} & susana.valencia@correounivalle.edu.co \hfill \textbf{Tel.:} (57) 3182155449 \\ \hline

\multicolumn{2}{|c|}{\textbf{Director}} \\ \hline
\textbf{Nombre:} & Humberto Loaiza Correa \hfill \textbf{Tel.:} (57) 6023212100 ext. 7654\\ \hline
\textbf{Correo electrónico:} & humberto.loaiza@correounivalle.edu.co \\ \hline
\textbf{Institución:} & Universidad del Valle - Percepción y sistemas inteligentes (PSI) \\ \hline

   



\end{tabular}
\end{table}
